\section{Schlussfolgerungen und Ausblick}\label{sec:conclusion}

Der anhaltende Trend steigender Motorisierung erhöht den Druck auf städtische Infrastrukturen durch vermeidbaren Suchverkehr kontinuierlich und rückt technologische Lösungsansätze wie \ac{iot}-basierte \acp{sps} und serverlose Cloud-Architekturen in den Mittelpunkt. Angesichts des wachsenden Wettbewerbsdrucks zur effizienteren Auslastung ihrer Stellplätze steht für Parkhausbetreiber die wirtschaftliche Sinnhaftigkeit im Vordergrund. Dabei ist entscheidend, inwieweit die notwendigen Investitionen in die Sensorik und Cloud-Infrastruktur eines \ac{sps} durch den zu erwartenden wirtschaftlichen Nutzen gerechtfertigt werden können. Folglich verlagert sich der Fokus von der technischen Umsetzbarkeit hin zur Ermittlung von Kennzahlen als Grundlage für Investitionsentscheidungen.

Um diese Lücke zu adressieren, stellt Drive \& Decide einen methodischen Ansatz vor, der auf einem prototypischen, \ac{iot}-basierten \ac{sps} mit einer serverlosen Architektur auf \ac{aws} aufbaut. Der Versuchsaufbau kombiniert ein 3D-gedrucktes Modell einer Parkgarage mit einer Cloud-Anbindung über \ac{mqtt}, \ac{aws} \ac{iot} Core, \ac{aws} Lambda und Amazon DynamoDB. Auf Basis dieses Prototyps wird ein Kosten-Nutzen-Modell aus Sicht eines Parkhausbetreibers entwickelt. Dieses stellt einmalige Investitionskosten (\ac{capex}) und laufende, nutzungsabhängige Cloud-Kosten (\ac{opex}) dem potenziellen Mehrerlös gegenüber.

Eine zentrale Aussage dieses Papers ist, dass Investitionsentscheidungen für ein \ac{sps} auf nachvollziehbaren Kennzahlen wie \ac{tco}, \ac{roi} und dem Break-Even-Punkt basieren sollten. Der durchgehend ereignisgetriebene Aufbau über serverlose Dienste ist dabei eine wesentliche Voraussetzung, um die laufenden Cloud-Kosten (\ac{opex}) belastbar zu erfassen und eine realistische Einschätzung der wirtschaftlichen Rentabilität zu ermöglichen.

Als Ausblick auf das Final Paper und die nächsten Arbeitsschritte wird der vorgestellte Versuchsaufbau physisch und virtuell umgesetzt. Das 3D-gedruckte Modell der Parkgarage wird mit Sensoren und Matchbox-Fahrzeugen ausgestattet, während parallel das \ac{aws}-Backend vollständig konfiguriert wird. Ziel ist es, anhand realer Messwerte und lastabhängiger Cloud-Kosten das entwickelte Kosten-Nutzen-Modell quantitativ zu evaluieren und Parkhausbetreibern eine empirisch fundierte Entscheidungsgrundlage zu liefern.
